\section{Environment Design and Implementation}
\label{ch:the_environment}

Before detailing the neuro-symbolic architectures developed in this thesis, this chapter establishes the experimental testbed. Testing an LLM-guided MARL system requires an environment with both physical complexity (movement, spatial reasoning) and logical dependencies (tasks gated by preconditions). To this end, we engineered a custom 2D multi-agent grid-world environment from scratch.

\subsection{Environment Design \& Mechanics}
\label{sec:env_design}

The environment operates as a cooperative spatial grid where multiple agents must jointly gather resources, craft tools, and defeat an adversary to achieve a final objective. Unlike standard dense-reward environments, where agents receive incremental gradient signals based on Euclidean distance to a target, this environment utilizes a logical sparse-reward formulation. The agents receive extrinsic rewards solely upon the completion of discrete logical milestones (subtasks).

The overarching philosophy of this design is mathematical austerity. The environment acts as a consistent judge, providing no incremental guidance or continuous signals. It only issues a constant time penalty to encourage speed and milestone bonuses for completing predefined physical tasks.

\subsubsection{Extrinsic Reward Structure}
The environment reward $r_t^{\text{env}}$ is the true, unmodified signal from the environment. It consists of a constant time-step penalty of $-0.01$ to encourage speed, plus milestone bonuses upon successfully completing a physical subtask:

\begin{table}[H]
\centering
\caption{Environment reward structure from \texttt{crafting\_env.py}.}
\label{tab:env_rewards}
\begin{tabular}{@{}lc@{}}
\toprule
\textbf{Event} & \textbf{Reward (shared)} \\
\midrule
Step penalty      & $-0.01$ \\
Collect Wood      & $+2.0$  \\
Collect Stone     & $+2.0$  \\
Craft Pickaxe     & $+3.0$  \\
Mine Iron         & $+2.0$  \\
Craft Sword/Armor & $+3.0$  \\
Build Bridge      & $+3.0$  \\
Defeat Enemy      & $+10.0$ \\
Mine Gold (win)   & $+15.0$ \\
\bottomrule
\end{tabular}
\end{table}

\subsubsection{The Dependency Graph (DAG)}
\label{subsec:env_dag}

The fundamental exploration challenge of the environment is formulated as a Directed Acyclic Graph (DAG) of dependencies. Agents cannot arbitrarily navigate to the final goal. They must coordinate to satisfy a rigid sequence of material and infrastructural preconditions. The required topological task progression is illustrated in Figure \ref{fig:env_dag}.

\begin{figure}[H]
    \centering
    \begin{tikzpicture}[
        node distance=1.5cm,
        box/.style={rectangle, draw=black, thick, fill=blue!10, text width=3cm, align=center, rounded corners, minimum height=1cm},
        gather/.style={box, fill=green!10},
        craft/.style={box, fill=orange!10},
        combat/.style={box, fill=red!10},
        arrow/.style={-{Stealth[scale=1.2]}, thick, draw=gray!80}
    ]
        % Row 1
        \node[gather] (wood) at (3, 0) {1a. Collect Wood (x2)};
        \node[gather] (stone) at (7, 0) {1b. Collect Stone};

        % Row 2
        \node[craft] (bridge) at (0, -2) {2a. Build Bridge};
        \node[craft] (pickaxe) at (5, -2) {2b. Craft Pickaxe};
        
        % Row 3
        \node[gather] (iron) at (5, -4) {3. Mine Iron (x2)};
        
        % Row 4
        \node[craft] (sword) at (3, -6) {4a. Craft Sword};
        \node[craft] (armor) at (7, -6) {4b. Craft Armor};

        % Row 5
        \node[combat] (enemy) at (3, -8.5) {5. Defeat Enemy};
        
        % Row 6
        \node[gather] (gold) at (3, -10.5) {6. Mine Gold};

        % Edges Row 1 -> 2
        \draw[arrow] (wood.west) -| (bridge.north);
        
        \coordinate (pickaxe_join) at (5, -1.2);
        \draw[thick, draw=gray!80, rounded corners=4pt] (wood.east) -| (pickaxe_join);
        \draw[thick, draw=gray!80, rounded corners=4pt] (stone.west) -| (pickaxe_join);
        \draw[arrow] (pickaxe_join) -- (pickaxe.north);

        % Edges Row 2 -> 3
        \draw[arrow] (pickaxe.south) -- (iron.north);
        
        % Edges Row 3 -> 4
        \coordinate (iron_split) at (5, -4.8);
        \draw[thick, draw=gray!80, rounded corners=4pt] (iron.south) -- (iron_split);
        \draw[arrow] (iron_split) -| (sword.north);
        \draw[arrow] (iron_split) -| (armor.north);
        
        % Edges to Enemy
        \draw[arrow] (bridge.south) |- (enemy.west);
        \draw[arrow] (sword.south) -- (enemy.north);
        \draw[arrow] (armor.south) |- (enemy.east);
        
        % Edge to Gold
        \draw[arrow] (enemy.south) -- (gold.north);

    \end{tikzpicture}
    \caption[Environment Dependency DAG]{The strict dependency Directed Acyclic Graph (DAG) governing environment progression. Agents must jointly acquire basic resources to clear logical bottlenecks. Notably, Wood is required for both the Pickaxe and the Bridge, while Iron supplies both Sword and Armor.}
    \label{fig:env_dag}
\end{figure}

As the topological depth of this chain increases, the probability of agents discovering the complete sequence through random exploration decays. This structure isolates the exploration bottleneck that our neuro-symbolic models are designed to address.

\subsubsection{Action and Observation Spaces}
\label{subsec:env_spaces}

\subsection{Partial Observability and the Action Space}
To simulate realistic constraints, the agents do not possess global vision of the grid. Instead, they operate under strict partial observability. At each time step $t$, agent $i$ receives a localized observation vector $o_t^{(i)} \in \mathbb{R}^{15}$ (or $\mathbb{R}^{25}$ in Chapter 6), which includes: their own $(x,y)$ coordinates and inventory status. The environment hides target locations (e.g., resources, workbenches) from the state vector. This requires agents to explore the spatial grid to locate objectives. To further prevent the policies from overfitting to a static sequence of movements, the agents' starting coordinates are randomized at the beginning of every episode.

The discrete action space consists of five parameters: four cardinal movement directions (Up, Down, Left, Right) and one semantic ``Interact'' action. The interaction mechanic is context-sensitive. For instance, triggering the action while adjacent to a forest tile chops wood, whereas executing it at a workbench (contingent upon possessing the requisite inventory) initiates tool crafting.

\subsubsection{Visual Representation}
\label{subsec:env_visual}

Figure~\ref{fig:env_screenshots} provides a visual representation of the discrete grid-world environment utilized throughout the MARL training process. The interface renders the underlying state matrix, allowing human observers to track the spatial coordinates of the autonomous agents. While the core simulation runs headlessly to maximize throughput, this renderer serves as a tool for visual debugging and policy interpretation. 

\begin{figure}[H]
    \centering
    \includegraphics[width=\textwidth]{../images/capture_from_the_game.png}
    \vspace{1ex}
    {\raggedleft \footnotesize Screenshot by \authorname~2026. Sprites generated by Gemini 1.5 Pro, Google. Prompt used: ``generate each sprites with different state if needed like enemy alive and dead and make sure we can distinguish each elements'' [generated July 10, 2026]. \par}
    \caption[Visual Rendering of the MARL Environment]{Visual rendering of the custom MARL environment. The interface explicitly renders the underlying state matrix, allowing human observers to track the spatial coordinates of the autonomous agents as they navigate the map.}
    \label{fig:env_screenshots}
\end{figure}

All visual assets utilized in this rendering engine were generated using the Gemini 1.5 Pro model via the Antigravity coding assistant platform.

\subsection{Architecture \& Technology Stack}
\label{sec:env_tech}

Engineering a custom MARL environment demands a careful selection of the underlying technology stack, requiring a balance between visual debugging capabilities and simulation throughput.

Why Not a Game Engine? Initially, open-source game engines such as Godot~\parencite{godot} were considered for the simulation backbone. Godot provides a visual editor, collision physics, and a node-based hierarchy, making it an option for traditional game development.

The final constraint in our environment design is the hardware profile. Because MARL requires processing thousands of physical state transitions per second to calculate gradients, introducing an LLM query directly into the training loop introduces significant computational overhead.

As a result, the environment was engineered natively in Python, utilizing NumPy~\parencite{numpy} as the core mathematical engine and Pygame~\parencite{pygame} as a lightweight visual wrapper. This matrix-first approach offers architectural advantages for reinforcement learning:
\begin{itemize}
    \item \textbf{Direct Tensor Integration:} The underlying grid, entities, and inventory states exist purely as NumPy matrices. This avoids the serialization and deserialization overhead typically encountered when passing state data between a Python-based RL training script and an external C++ game engine.
    \item \textbf{Seamless NumPy-to-Game Translation:} Because the environment state is natively represented as a matrix, Pygame can translate these numerical coordinates into a visual interface for human debugging. The mathematical logic (kinematics, collisions) is decoupled from the rendering loop.
    \item \textbf{Headless Vectorization:} During the actual training loop, the Pygame visual renderer is disabled. The PyTorch models interact directly with the raw NumPy arrays, allowing a vectorized wrapper to step thousands of environments simultaneously in parallel.
\end{itemize}

\subsection{Libraries \& Infrastructure}
\label{sec:env_libs}

The codebase relies on a modern, standardized stack of Python libraries to seamlessly interface the custom environment with the deep neural networks and the large language models:

\begin{itemize}
    \item \textbf{MARL Interface:} The environment conforms to the PettingZoo API \parencite{Terry_Black_Grammel_Jayakumar_Hari_Sullivan_Santos_Perez_Horsch_Dieffendahl_et_al._2021}, which has established itself as the standard for multi-agent RL (analogous to the Gymnasium API for single-agent RL). Utilizing the Parallel API variant ensures the environment integrates natively with synchronous MAPPO tensor rollouts.
    \item \textbf{Deep Learning Framework:} PyTorch \parencite{Paszke_Gross_Massa_Lerer_Bradbury_Chanan_Killeen_Lin_Gimelshein_Antiga_et_al_2019} serves as the computational backbone for designing, compiling, and optimizing the MAPPO Actor and Critic networks, leveraging GPU acceleration for Generalized Advantage Estimation (GAE) and gradient descent.
    \item \textbf{LLM Integration:} To interface with the semantic Meta-Controllers, the Transformers library by Hugging Face is deployed \parencite{Wolf_Debut_Sanh_Chaumond_Delangue_Moi_Cistac_Rault_Louf_Funtowicz_et_al._2020}. For the Hierarchical Reinforcement Learning phase (which necessitates rapid, localized inference), the language models are quantized using bitsandbytes (QLoRA) \parencite{Dettmers_Pagnoni_Holtzman_Zettlemoyer_2023} and served via the Ollama framework~\parencite{ollama_2024} to strictly bound the latency introduced during the training loop.
    \item \textbf{Experiment Tracking:} To monitor the highly stochastic training process, TensorBoard is integrated directly into the PyTorch training loop. This allows for real-time visualization of critical metrics such as the episodic return, actor-critic loss curves, and policy entropy across millions of environment transitions.
\end{itemize}

\subsection{Hardware and Inference Architecture}
\label{sec:env_hardware}

A critical bottleneck in neuro-symbolic MARL is the severe latency introduced by the LLM during the training loop. This thesis relies on two distinct inference paradigms, online cloud APIs and localized serving, dictated by the constraints of the specific algorithms.

For the Dynamic Reward Shaping architecture, which requires infrequent, highly complex offline generation, this study leveraged the Gemini 2.5 Flash model hosted via Google AI Studio\footnote{https://aistudio.google.com/}. However, integrating cloud-based models directly into the high-frequency Hierarchical Reinforcement Learning (HRL) training loop proved intractable. The MAPPO environment requires thousands of synchronous environment steps per episode. Utilizing an online model for the HRL orchestrator generated an excessive volume of requests, rapidly exhausting API rate limits and introducing network latency that stalled the GPU-bound RL rollouts.

To overcome this, the HRL training loop was executed entirely locally on a consumer-grade workstation equipped with an NVIDIA RTX 5070 GPU (12GB VRAM). Running a foundation model locally alongside the PyTorch MAPPO actor-critic tensors within a strict 12GB memory envelope necessitated specific architectural compromises. To address this, the Ollama framework~\parencite{ollama_2024} was deployed to provide highly optimized, low-latency local inference.

As detailed in Section~\ref{sec:qwen_architecture}, the Qwen 2.5 7B architecture was selected specifically to satisfy this strict 12GB memory envelope. When combined with 4-bit quantization via Ollama, it fits comfortably alongside the MARL experience replay buffers without displacing the PyTorch actor-critic tensors.